\chapter{Experiment 1: Fixed-Point Tracking with Rotated Stiffness Directions}
\label{app:exp1}

\section{Overview}

This appendix provides the implementation details for
Experiment~1. The experiment tracks a single fixed Cartesian point using a full
base-frame stiffness matrix. The matrix is built from a diagonal local stiffness
matrix and a rotated local stiffness frame:

\begin{equation}
\mathbf{K}_{p,\mathrm{base}}
=
\mathbf{W}_{\mathrm{local}\rightarrow\mathrm{base}}
\mathbf{K}_{p,\mathrm{local}}
\mathbf{W}_{\mathrm{local}\rightarrow\mathrm{base}}^T.
\end{equation}

The purpose is to demonstrate that full stiffness and damping matrices are not
only relevant for surface or wall tasks. They also appear naturally in fixed
point tracking when the desired principal stiffness directions are rotated
relative to the base frame.

\section{C++ Implementation}

The listing below shows the core of the controller. The same indexing style is
used throughout the implementation: matrices with the suffix
\texttt{\_local} are defined in the local stiffness frame; matrices with the
suffix \texttt{\_base} are used in the controller.

\begin{lstlisting}[style=cppstyle,
  caption={Experiment 1: fixed-point tracking with rotated stiffness directions},
  label={lst:exp1_rotated_cpp}]
#include <Eigen/Dense>
#include <franka/model.h>
#include <franka/robot.h>
#include <cmath>

// Desired fixed point
Eigen::Vector3d p_des(0.50, 0.00, 0.40);

// Rotation of local stiffness frame relative to base frame
double alpha = 20.0 * M_PI / 180.0;  // roll
double beta  = 30.0 * M_PI / 180.0;  // pitch
double gamma = 40.0 * M_PI / 180.0;  // yaw

Eigen::AngleAxisd Rx(alpha, Eigen::Vector3d::UnitX());
Eigen::AngleAxisd Ry(beta,  Eigen::Vector3d::UnitY());
Eigen::AngleAxisd Rz(gamma, Eigen::Vector3d::UnitZ());

Eigen::Matrix3d W_local_to_base =
    (Rz * Ry * Rx).toRotationMatrix();

// Local stiffness and damping
Eigen::Matrix3d Kp_local;
Kp_local << 1000.0, 0.0,   0.0,
            0.0,    300.0, 0.0,
            0.0,    0.0,   100.0;

Eigen::Matrix3d Dp_local;
Dp_local << 63.25, 0.0, 0.0,
            0.0, 34.64, 0.0,
            0.0, 0.0, 20.00;

// Transform local matrices into base frame
Eigen::Matrix3d Kp_base =
    W_local_to_base * Kp_local * W_local_to_base.transpose();

Eigen::Matrix3d Dp_base =
    W_local_to_base * Dp_local * W_local_to_base.transpose();

// Inside the 1 ms torque callback:
// p_EE: current end-effector position
// pdot: current end-effector translational velocity
// J: 6x7 geometric Jacobian
// tau_g, tau_c: gravity and Coriolis compensation

Eigen::Vector3d e_p = p_des - p_EE;
Eigen::Vector3d pdot_d = Eigen::Vector3d::Zero();
Eigen::Vector3d f = Kp_base * e_p + Dp_base * (pdot_d - pdot);

Eigen::Matrix<double,6,1> F;
F.head<3>() = f;
F.tail<3>().setZero();

Eigen::Matrix<double,7,1> tau_task = J.transpose() * F;
Eigen::Matrix<double,7,1> tau_null = Eigen::Matrix<double,7,1>::Zero();
Eigen::Matrix<double,7,1> tau_cmd  =
    tau_task + tau_null + tau_g + tau_c;
\end{lstlisting}

\section{Expected Numerical Values}

Using the angles
\(\alpha=20^\circ\), \(\beta=30^\circ\), and \(\gamma=40^\circ\), the
base-frame stiffness matrix is approximately

\begin{equation}
\mathbf{K}_{p,\mathrm{base}}
\approx
\begin{bmatrix}
605.1 & 359.0 & -205.6\\
359.0 & 500.2 & 138.9\\
-205.6 & 138.9 & 294.7
\end{bmatrix}
\;\mathrm{N/m}.
\end{equation}

For an initial error

\begin{equation}
\mathbf{e}_p(0)=
\begin{bmatrix}
0.01\\0\\0
\end{bmatrix}
\mathrm{m},
\end{equation}

and zero initial velocity, the expected initial force is

\begin{equation}
\mathbf{f}(0)
=
\mathbf{K}_{p,\mathrm{base}}\mathbf{e}_p(0)
\approx
\begin{bmatrix}
6.05\\3.59\\-2.06
\end{bmatrix}
\mathrm{N}.
\end{equation}

This confirms that a pure \(x\)-error in the base frame does not necessarily
lead to a pure \(x\)-force. The force follows the rotated stiffness field.

\section{Python Verification Code}

The following Python script verifies the matrix construction and force
calculation.

\begin{lstlisting}[style=cppstyle, language=Python,
  caption={Python verification of the rotated stiffness matrix},
  label={lst:exp1_rotated_python}]
import numpy as np

def Rx(a):
    c, s = np.cos(a), np.sin(a)
    return np.array([[1,0,0],[0,c,-s],[0,s,c]])

def Ry(b):
    c, s = np.cos(b), np.sin(b)
    return np.array([[c,0,s],[0,1,0],[-s,0,c]])

def Rz(g):
    c, s = np.cos(g), np.sin(g)
    return np.array([[c,-s,0],[s,c,0],[0,0,1]])

alpha = np.deg2rad(20.0)
beta  = np.deg2rad(30.0)
gamma = np.deg2rad(40.0)

W = Rz(gamma) @ Ry(beta) @ Rx(alpha)

K_local = np.diag([1000.0, 300.0, 100.0])
K_base = W @ K_local @ W.T

e = np.array([0.01, 0.0, 0.0])
f = K_base @ e

print("K_base =")
print(np.round(K_base, 2))
print("f =", np.round(f, 2), "N")
\end{lstlisting}

\section{Parameter Table}

\begin{center}
\renewcommand{\arraystretch}{1.3}
\begin{tabular}{@{}llll@{}}
\toprule
Parameter & Symbol & Value & Unit \\
\midrule
Desired point & \(\mathbf{p}_d\) & \([0.50,0.00,0.40]^T\) & m \\
Initial error & \(\mathbf{e}_p(0)\) & \([0.01,0,0]^T\) & m \\
Roll angle & \(\alpha\) & 20 & deg \\
Pitch angle & \(\beta\) & 30 & deg \\
Yaw angle & \(\gamma\) & 40 & deg \\
Local stiffness & \(\mathbf{K}_{p,\mathrm{local}}\) & \(\operatorname{diag}(1000,300,100)\) & N/m \\
Virtual mass & \(\mathbf{M}_p\) & \(\mathbf{I}_3\) & kg \\
Control frequency & \(f_s\) & 1000 & Hz \\
\bottomrule
\end{tabular}
\end{center}

\section{Relation to the Main Thesis}

This experiment is intentionally simple: the desired point is fixed, and no
external force is applied. The important point is not the position itself but
the force direction generated by the full stiffness matrix. It demonstrates
that the desired point defines the centre of the virtual spring, whereas
\(\mathbf{K}_{p,\mathrm{base}}\) defines the shape and orientation of the
spring behaviour around that point.
